\section{The comparison-cohomology problems recovered in cote 18}

\subsection{Original question and precise modern scope}
The original source is Grothendieck's Montpellier archive, cote 18,
physical PDF page 43, manuscript and printed footer 42, headed
``Pbs.''  Its four numbered headings read ``Betti, De-Rham,''
``Betti, $\ell$-adique,'' ``DeRham, $p$-adique,'' and
``D\'efinition de cohomologie $p$-adique en car $p>0$?''
The first heading also branches into comparison problems involving
different embeddings and additional structures.  We separate those
branches from the four precise cohomological assertions below.
The original sheet was visually inspected in this run.

These are answers by known comparison theorems, with their hypotheses
retained.  They are not findings newly proved from RH.  The IDs
\texttt{KN-COMP01--04} identify source-checked answers, not claims of
new comparison theorems.  The unrelated uppercase $K$ on physical
page 45/footer 44 is still not defined by the recovered adjacent
sheets; no meaning is assigned to it here.

\paragraph{KN-COMP01: Betti--de Rham comparison.}
Let $X$ be a smooth separated finite-type scheme over $\mathbf C$.
There is a canonical comparison
\[
 \mathbb H^i(X,\Omega^\bullet_{X/\mathbf C})
       \simeq H^i(X^{\rm an},\mathbf C).
 \tag{COMP.1}
\]
This is Grothendieck (1966), Theorem $1'$,
\emph{Publications Math\'ematiques de l'IH\'ES} 29, p.~96;
Theorem 1, p.~95, gives the smooth affine form.  The theorem is
cited, not reproved.  If a variety is defined over a subfield and
an embedding into $\mathbf C$ is specified, (COMP.1) is applied to
that complex base change.  This answers the cohomological comparison
for each specified embedding.  A preferred rational identification
between the Betti fibre functors of two different embeddings is
additional data; the complex comparison alone does not specify it.

\paragraph{KN-COMP02: Betti--\texorpdfstring{$\ell$}{ell}-adic comparison.}
For a separated finite-type complex scheme $X$ and every finite
constant abelian coefficient group $A$, the analytification comparison gives
\[
 H^i_{\rm et}(X,A)\simeq H^i(X^{\rm an},A).
 \tag{COMP.2}
\]
The precise general input is Artin, SGA~4, Expos\'e XVI,
Theorem 4.1; in the inspected re-edition its statement occupies
pp.~13--14.  For $f:X\to S$, it requires $f$ proper or the
coefficient sheaf constructible, and identifies the corresponding
higher direct images.  Constant finite $A$ is constructible, so the
absolute comparison above is included.  The familiar $\mathbf Z_\ell$
and $\mathbf Q_\ell$ comparisons use the compatible finite-level maps.
For example, for smooth proper $X$ the finite groups at each level
satisfy the Mittag--Leffler condition: for fixed level the descending
images stabilize.  Taking the inverse limit in (COMP.2), then
inverting $\ell$, gives the stated adic comparisons.  The comparison
is over the chosen complex base change; it does not supply an
unspecified rational period-torsor point.

\paragraph{KN-COMP03: crystalline cohomology in characteristic
\texorpdfstring{$p$}{p}.}
For an $\mathbf F_p$-scheme $S$, the crystalline site relative to
$\mathbf Z_p$ uses nilpotent divided-power thickenings of opens of
$S$, compatible with the canonical divided powers on $(p)\subset\mathbf Z_p$; its structure sheaf assigns the functions on the thickening.
Its derived global sections define $R\Gamma_{\rm crys}(S/\mathbf Z_p)$.
A crystal has the prescribed base-change isomorphisms along these
thickenings.  This is the directly inspected definition in
Bhatt--de Jong (2011), \S1, pp.~1--2.
Their Theorem 3.6, p.~8, computes crystalline cohomology by the
de Rham complex on a completed divided-power envelope.  Their
Corollary 3.8, p.~9, gives, for $X$ proper smooth over $\mathbf Z_p$
and $S=X\otimes\mathbf F_p$,
\[
 R\Gamma_{\rm crys}(S/\mathbf Z_p)
        \simeq R\Gamma(X,\Omega^\bullet_{X/\mathbf Z_p}).
 \tag{COMP.3}
\]
The left side is intrinsic to $S$, so the comparison also explains
the lift-independent Frobenius on this de Rham cohomology.
These are known results, with no assumption that the special fibre
has a chosen lift in the definition itself.

\paragraph{KN-COMP04: de Rham--\texorpdfstring{$p$}{p}-adic and
integral mixed-characteristic comparisons.}
Let $K/\mathbf Q_p$ be complete discretely valued with perfect
residue field, $C$ a complete algebraically closed extension, and
$X_0/K$ proper smooth.  Bhatt--Morrow--Scholze (2018), Theorem 1.7,
author version p.~5, gives the canonical comparison
\[
 H^i_{\rm dR}(X_0/K)\otimes_K B_{\rm dR}
    \simeq H^i_{\rm et}(X_0\widehat\otimes_K C,\mathbf Z_p)
                    \otimes_{\mathbf Z_p}B_{\rm dR}.
 \tag{COMP.4}
\]
For a proper smooth formal scheme $\mathfrak X/\mathcal O_C$, their
Theorem 1.8, pp.~5--6, supplies the perfect complex
$A=R\Gamma_{A_{\rm inf}}(\mathfrak X)$, where
$A_{\rm inf}=W(\mathcal O_C^\flat)$, with Frobenius and comparisons
\[
\begin{aligned}
 A\otimes^{L}_{A_{\rm inf}}W(\kappa_C)
   &\simeq R\Gamma_{\rm crys}(\mathfrak X_{\kappa_C}/W(\kappa_C)),\\
 A\otimes^{L}_{A_{\rm inf}}\mathcal O_C
   &\simeq R\Gamma_{\rm dR}(\mathfrak X),\\
 A[1/\mu]&\simeq
 R\Gamma_{\rm et}(\mathfrak X_C,\mathbf Z_p)
       \otimes^L_{\mathbf Z_p}A_{\rm inf}[1/\mu].
\end{aligned}
 \tag{COMP.5}
\]
Here $\kappa_C$ is the residue field and
$\mu=[\epsilon]-1$ comes from compatible $p$-power roots of unity.
These derived comparisons retain torsion information; replacing them
by unqualified integral cohomology-group isomorphisms would lose
their hypotheses.  The cited theorem also treats the crystalline
$A_{\rm crys}$ specialization.

\subsection{Consequences for the requested bridge}
The four numbered headings therefore have precise affirmative
cohomological answers in the scopes displayed above.  They explain
how Betti, de Rham, crystalline and \emph{generic-fibre} adic
realizations can belong to one arithmetic-geometric comparison
system.  The period rings and the divided-power or integral
comparison complexes are part of that system's mathematical data.
The adopted RH and Dirichlet-GRH theorems determine the zero locations
of the abelian Artin--Tate factors already isolated in BR-AT-01.
The comparison results identify realizations; they do not by
themselves construct the algebraic cross-correspondences still
required in CC-XC01--04, or choose a rational point on the period
torsors in the remaining cote 18 questions.

\paragraph{Primary references and inspected locators.}
Grothendieck, A. (1966). On the de Rham cohomology of algebraic
varieties. \emph{Publications Math\'ematiques de l'IH\'ES, 29},
95--103. \url{https://doi.org/10.1007/BF02684807}.
Theorems 1 and $1'$, pp.~95--96; primary scan:
\url{https://www.numdam.org/article/PMIHES_1966__29__95_0.pdf}.

Artin, M. (1973). Th\'eor\`eme de changement de base par un
morphisme lisse, et applications. In M. Artin, A. Grothendieck,
\& J.-L. Verdier (Eds.), \emph{Th\'eorie des topos et cohomologie
\'etale des sch\'emas (SGA 4)}, Tome 3, Expos\'e XVI,
Theorem 4.1. Lecture Notes in Mathematics 305. Springer.
Inspected re-edition, version 71766d9 (30 July 2024), pp.~13--14:
\url{https://www.normalesup.org/~forgogozo/SGA4/16/16.pdf}.
The version-specific pages are not asserted as original pagination.

Bhatt, B., \& de Jong, A. J. (2011).
\emph{Crystalline cohomology and de Rham cohomology}
[Preprint, arXiv:1110.5001]. \S1 pp.~1--2; Theorem 3.6 p.~8;
Corollary 3.8 p.~9.
\url{https://www.math.columbia.edu/~dejong/papers/crystalline-comparison.pdf}.

Bhatt, B., Morrow, M., \& Scholze, P. (2018).
Integral $p$-adic Hodge theory.
\emph{Publications Math\'ematiques de l'IH\'ES, 128}, 219--397.
\url{https://doi.org/10.1007/s10240-019-00102-z}.
Inspected author version dated 8 January 2019, Theorems 1.7--1.8,
pp.~5--6:
\url{https://www.math.uni-bonn.de/people/scholze/integralpadicHodge.pdf}.
